% Part: first-order-logic
% Chapter: completeness
% Section: completeness-theorem

\documentclass[../../../include/open-logic-section]{subfiles}

\begin{document}

\iftag{FOL}
      {\olfileid{fol}{com}{cth}}
      {\olfileid{pl}{com}{cth}}

% SEGMENT OLP-0134-S01
\olsection{நிறைவு தேற்றம்}

% SEGMENT OLP-0134-S02
\begin{explain}
நமது முடிவுகளை ஒன்றிணைப்போம்: நிறைவு தேற்றத்தை அடைகிறோம்.
\end{explain}

% SEGMENT OLP-0134-S03
\begin{thm}[நிறைவு தேற்றம்]
\ollabel{thm:completeness}
$\Gamma$ என்பது !!{sentence}s கொண்ட ஒரு கணம் என்க. $\Gamma$ முரணற்றது எனில்,
அது நிறைவுசெய்யத்தக்கது.
\end{thm}

\begin{proof}
$\Gamma$ முரணற்றது என்க.\iftag{FOL}{ \olref[hen]{lem:henkin} படி,
$\Gamma' \supseteq \Gamma$ ஆக இருக்கும் செறிவுற்ற முரணற்ற கணம் ஒன்று
உள்ளது.}{} \olref[lin]{lem:lindenbaum} படி,
$\Gamma^* \supseteq \iftag{FOL}{\Gamma'}{\Gamma}$ ஆகவும் முரணற்றதாகவும்
!!{complete} ஆகவும் இருக்கும் கணம் ஒன்று உள்ளது.
\iftag{FOL}{
  $\Gamma' \subseteq \Gamma^*$ என்பதால், ஒவ்வொரு !!{formula}~$!A(x)$ க்கும்
  பின்வரும் வடிவிலான !!a{sentence} ஐ~$\Gamma^*$ கொண்டுள்ளது:
  \iftag{prvEx}
        {$\lexists[x][!A(x)] \lif !A(c)$}
        {$\lnot\lforall[x][!A(x)] \lif \lnot !A(c)$}.
  ஆகவே~$\Gamma^*$ செறிவுற்றது. $\Gamma$ ஆனது~$\eq$ ஐக் கொண்டிருக்கவில்லை
  எனில்,}{அப்போது} \olref[mod]{lem:truth} படி,
$\iftag{FOL}{\Sat{M(\Gamma^*)}}{\pSat{v(\Gamma^*)}}{!A}$ அப்போதும் அப்போது
மட்டுமே $!A \in \Gamma^*$. குறிப்பாக, எல்லா $!A \in \Gamma$ களுக்கும்
$\iftag{FOL}{\Sat{M(\Gamma^*)}}{\pSat{v(\Gamma^*)}}{!A}$ என்பது இதிலிருந்து
கிடைக்கிறது; ஆகவே~$\Gamma$ நிறைவுசெய்யத்தக்கது.\iftag{FOL}{
  $\Gamma$ ஆனது~$\eq$ ஐக் கொண்டிருந்தால், \olref[ide]{lem:truth} படி எல்லா
  !!{sentence}s~$!A$ க்கும் $\Sat{\equivclass{M}{\approx}}{!A}$ அப்போதும்
  அப்போது மட்டுமே $!A \in \Gamma^*$. குறிப்பாக, எல்லா $!A \in \Gamma$
  களுக்கும் $\Sat{\equivclass{M}{\approx}}{!A}$; ஆகவே~$\Gamma$
  நிறைவுசெய்யத்தக்கது.}{}
\end{proof}

% SEGMENT OLP-0134-S04
\begin{cor}[நிறைவு தேற்றம், இரண்டாம் வடிவம்]
\ollabel{cor:completeness}
எல்லா~$\Gamma$ களுக்கும் !!{sentence}s~$!A$ க்கும்: $\Gamma \Entails !A$
எனில் $\Gamma \Proves !A$.
\end{cor}

% SEGMENT OLP-0134-S05
\begin{proof}
\olref{cor:completeness} மற்றும் \olref{thm:completeness} ஆகியவற்றிலுள்ள
$\Gamma$ கள் அனைத்தளவையிடப்பட்டுள்ளன என்பதைக் கவனிக்க. குழப்பத்தைத்
தவிர்ப்பதற்காக, \olref{thm:completeness} ஐ வேறொரு மாறியைப் பயன்படுத்தி
மீண்டும் கூறுவோம்: !!{sentence}s கொண்ட எந்தக் கணம்~$\Delta$ க்கும்,
$\Delta$ முரணற்றது எனில் அது நிறைவுசெய்யத்தக்கது. எதிர்நிலைமாற்றத்தால்,
$\Delta$ நிறைவுசெய்யத்தக்கதல்ல எனில், $\Delta$ முரணுடையது. இதைப் பயன்படுத்தி
முடிவுரையை நிறுவுவோம்.

$\Gamma \Entails !A$ என்க. அப்போது \olref[syn][sem]{prop:entails-unsat}
படி, $\Gamma \cup \{\lnot !A\}$ நிறைவுசெய்யத்தக்கதல்ல. $\Gamma \cup
\{\lnot !A\}$ ஐ நமது~$\Delta$ ஆக எடுத்துக்கொண்டால்,
\olref{thm:completeness} இன் முந்தைய வடிவம் $\Gamma \cup \{\lnot !A\}$
முரணுடையது எனத் தருகிறது. பின்வரும் மேற்கோளின்படி
\iftag{FOL}{%
  \tagrefs{prfAX/{fol:axd:prv:prop:prov-incons},
    prfSC/{fol:seq:prv:prop:prov-incons},
    prfND/{fol:ntd:prv:prop:prov-incons},
    prfTab/{fol:tab:prv:prop:prov-incons}}}{%
  \tagrefs{prfAX/{pl:axd:prv:prop:prov-incons},
    prfSC/{pl:seq:prv:prop:prov-incons},
    prfND/{pl:ntd:prv:prop:prov-incons},
    prfTab/{pl:tab:prv:prop:prov-incons}}},
$\Gamma \Proves !A$.
\end{proof}

% SEGMENT OLP-0134-S06
\tagprob{FOL}
\begin{prob}
\olref[fol][com][cth]{cor:completeness} ஐப் பயன்படுத்தி
\olref[fol][com][cth]{thm:completeness} ஐ நிறுவுக; இதன்மூலம் நிறைவு
தேற்றத்தின் இரு வடிவங்களும் சமானமானவை எனக் காட்டுக.
\end{prob}
\tagendprob

% SEGMENT OLP-0134-S07
\tagprob{notFOL}
\begin{prob}
\olref[pl][com][cth]{cor:completeness} ஐப் பயன்படுத்தி
\olref[pl][com][cth]{thm:completeness} ஐ நிறுவுக; இதன்மூலம் நிறைவு
தேற்றத்தின் இரு வடிவங்களும் சமானமானவை எனக் காட்டுக.
\end{prob}
\tagendprob

% SEGMENT OLP-0134-S08
\tagprob{FOL}
\begin{prob}
!!a{derivation} முறைமை நிறைவானதாக இருப்பதற்கு, நிறைவுசெய்யத்தக்கதல்லாத
ஒவ்வொரு கணமும் முரணுடையது என நிறுவுமளவுக்கு அதன் விதிகள் வலுவாக இருக்க
வேண்டும். நிறைவை நிறுவுவதற்கு !!{derivation} விதிகளில் எவை தேவைப்பட்டன?
அவ்விதிகளில் ஏதேனும் நிறுவலில் எங்கும் பயன்படுத்தப்படவில்லையா? இக்கேள்விகளுக்கு
விடையளிக்க, \olref[fol][com][cth]{thm:completeness} இன் நிறுவலுக்கு இட்டுச்
செல்லும் எந்த முடிவுகளில் எந்த !!{derivation} விதிகள் பயன்படுத்தப்பட்டன
என்பதைக் காட்டும் பட்டியல் அல்லது வரைபடத்தை உருவாக்குக. இந்நிறுவல்களில்
வெளிப்படையாகக் கூறப்படாத விதிப் பயன்பாடுகளையும் குறிப்பிடுக.
\end{prob}
\tagendprob

% SEGMENT OLP-0134-S09
\tagprob{notFOL}
\begin{prob}
!!a{derivation} முறைமை நிறைவானதாக இருப்பதற்கு, நிறைவுசெய்யத்தக்கதல்லாத
ஒவ்வொரு கணமும் முரணுடையது என நிறுவுமளவுக்கு அதன் விதிகள் வலுவாக இருக்க
வேண்டும். நிறைவை நிறுவுவதற்கு !!{derivation} விதிகளில் எவை தேவைப்பட்டன?
அவ்விதிகளில் ஏதேனும் நிறுவலில் எங்கும் பயன்படுத்தப்படவில்லையா? இக்கேள்விகளுக்கு
விடையளிக்க, \olref[pl][com][cth]{thm:completeness} இன் நிறுவலுக்கு இட்டுச்
செல்லும் எந்த முடிவுகளில் எந்த !!{derivation} விதிகள் பயன்படுத்தப்பட்டன
என்பதைக் காட்டும் பட்டியல் அல்லது வரைபடத்தை உருவாக்குக. இந்நிறுவல்களில்
வெளிப்படையாகக் கூறப்படாத விதிப் பயன்பாடுகளையும் குறிப்பிடுக.
\end{prob}
\tagendprob

\end{document}
